\documentclass[a4paper, 11pt]{article}

% -- Coments
\usepackage{verbatim}

% ----- Fonts -----
% -- Fuente --
% \usepackage{fontspec}
% \setmonofont{JetBrainsMono Nerd Font}

\usepackage{babel}
\babelprovide[import, main]{spanish}

% -- Color --
\usepackage{xcolor}
%\definecolor{azul}{RGB}{00,33,99}
\definecolor{azul}{RGB}{35,72,180}

% -- Page Margin --
\usepackage[margin=1in]{geometry}

% -- Espaciados --
\newcommand{\Pspace}{0.5cm}
\newcommand{\Aspace}{0.2cm}

% -- Columnas --
\usepackage{multicol}

% -- Imagenes --
\usepackage{graphicx}
\usepackage{float}

% -- Matemáticas --
\usepackage{amsmath, amssymb}

% -- Gráficas --
\usepackage{pgfplots}
\pgfplotsset{compat=1.18}


\title
{
    Contexto Regional 2026-2 \\
    Golondrinas errantes. Mujeres sonorenses en la Revolución Mexicana - Vásquez Montaño, Rosario Margarita \\

    Profesor José Moreno Vega
}

\begin{document}

\begin{otherlanguage}{spanish}
\maketitle
\end{otherlanguage}

\begin{center}
    \begin{tabular}{r|l}
        \textbf{Expediente} & \textbf{Nombre} \\ \hline
        219208106 & Bórquez Guerrero Angel Fernando \\
    \end{tabular}
\end{center}

{\Large 1. ¿Cuál es el argumento central del autor?} \\
{\color{azul}
    La autora dice que la historia de la Revolución en Sonora se ha contando casi siempre desde los caudillos, los militares y los rebeldes, y que las mujeres quedaron en el silencio aunque participaron mucho. Según ella, las sonorenses fueron enfermeras, espías, contrabandistas, propagandistas y hasta combatientes. Las llama "golondrinas errantes" porque se movían por todo el estado y cruzaban la frontera. Al hacerlo rompieron con los papeles que la sociedad les tenía asignados y demostraron que podían actuar por cuenta propia. Por eso concluye que ya no se puede estudiar la Revolución sin incluirlas.
}
\vspace{1cm}

{\Large 2. ¿Qué evidencias utiliza para respaldarlo?} \par
{\color{azul}
Usa sobre todo documentos de archivo, como los expedientes de veteranas de la Revolución de la Secretaría de la Defensa Nacional, donde mujeres como Francisca Celis, Ernestina Agüero y Laura Rochin piden ser reconocidas por lo que hicieron. También usa memorias y testimonios, como La Rebelde de Leonor Villegas de Magnón y los recuerdos de Nicolás Bernal. Se apoya además en el periódico Regeneración, del movimiento magonista. Presenta casos concretos, como el de María de Jesús González, que se vistió de soldado, y el de Margartia Ortega, que fue torturada y fusilada en 1913. Por último, cita estudios de hisotria regional y de historia de las mujeres.
}
\vspace{1cm}

{\Large 3. ¿Qué relación hay con la lectura anterior?} \par
{\color{azul}
Las dos lecturas rescatan a gente que casi no aparece en la historia oficial: mujeres en la lectura 7 y los yaquis en la lectura 6. En ambos casos las autoras buscan a esas personas "escondidas" en archivos judiciales o militares, porque no hay muchos otros registros.
}

\end{document}
